% ==============================================================================
% content.tex — 自由意思とフレーム慣性 (日本語教科書スタイル本文)
% ==============================================================================

\section{序文}\label{sec:intro}

私たちは日常生活で、あるパラドックスによく直面します：
学習、運動、あるいは長期的に有益な活動を行う\emph{べきであると明確に分かっている}にもかかわらず、
その一歩を踏み出すことができないという経験です。
あたかも目に見えない「渦」が私たちを現在の状態（スマートフォンを眺める、横になる、ぼーっとする）に
吸い寄せているかのようです。
理性は「動け！」と叫んでいますが、身体と感情は鉛のように重く、現在の状態に留まろうとします。

これは単なる「怠慢」ではなく、構造化された分析が可能な現象です。
本稿では、この問題を「フレーム慣性理論 (Framework Inertia Theory)」という簡易的な形式枠組みを用いて記述し、
自由意思が最大活用されるタイミングを特定し、なぜ\emph{計画}が単なる選択肢ではなく\emph{必要不可欠}であるのかを導き出します。

% ─────────────────────────────────────────────────────────────────────────────
\section{フレーム慣性}\label{sec:inertia}

\begin{definition}[行動状態]\label{def:state}
\textbf{行動状態} $s \in \mathcal{S}$ とは、特定の活動パターンと、それに伴う物理的環境、認知的文脈、
および神経化学的な報酬プロファイルが組み合わさった持続的な状態を指す。
\end{definition}

\begin{examplebox}
ベッドに横になって短い動画を視聴し続けることは、行動状態 $s_0$ を構成します。
その特徴は：(i) 仰向けという物理的姿勢、(ii) 最小限の心理的努力で済む認知的背景、
(iii) 絶え間なく供給されるドーパミンの微小報酬です。
一方、机に向かって学習することは、高い認知負荷と報酬の遅延を伴う行動状態 $s_1$ です。
\end{examplebox}

\begin{definition}[フレーム慣性]\label{def:inertia}
ある行動状態 $s$ の\textbf{フレーム慣性} $I(s)$ とは、その状態 $s$ から他の状態へ遷移しようとする際に
直面する抵抗の大きさを指す。形式的には、
\[
  I(s) \;=\; \alpha \, R(s) \;+\; \beta \, C(s) \;+\; \gamma \, E(s),
\]
ここで、
\begin{itemize}[nosep]
  \item $R(s) \geq 0$ は現在の活動の\emph{瞬時報酬率}（神経的要素）、
  \item $C(s) \geq 0$ は\emph{認知的切り替えコスト}（認知的要素）、
  \item $E(s) \geq 0$ は\emph{環境的強化}（環境的要素）、
\end{itemize}
であり、$\alpha, \beta, \gamma > 0$ は個人の特性に基づく重み係数とする。
\end{definition}

\subsection{慣性の渦モデル}

物理学のポテンシャル関数のアナロジーを用いると、行動状態 $s$ は「勢阱（ポテンシャルの井戸）」として解釈できます。
状態 $s_0$ から $s_1$ へ遷移するためには、\textbf{活性化エネルギー障壁} $\Delta V$ を克服する必要があります。

\begin{theorem}[自由意思の最大レバレッジ]\label{thm:leverage}
自由意思は、常に一定の力で働くのではなく、\textbf{意思決定の窓（デシジョン・ウィンドウ）}---すなわち、フレーム慣性が局所的に最小となる瞬間において最大の効力を発揮する。
\end{theorem}

\begin{discussion}
意思決定の窓は、ポテンシャル障壁が低下する瞬間に現れます：
\begin{itemize}[nosep]
  \item \textbf{起床直後：} 前日の行動フレームがまだ完全にロードされていないため、低コストで新しいフレームを挿入可能。
  \item \textbf{活動の区切り：} タスク完了の瞬間、慣性が一時的にゼロになる「隙」が生まれる。
  \item \textbf{環境の遷移：} 図書館に入る、あるいはジムの入り口に到着するなどの物理的環境の変化は、古いフレームを強制的に破棄する。
\end{itemize}
したがって、最善の戦略は「慣性と向き合って力技で戦い続ける」ことではなく、
「\emph{慣性が最も弱まった瞬間に意思決定を行い、その後の持続は新しい慣性に任せる}」ことです。
\end{discussion}

% ─────────────────────────────────────────────────────────────────────────────
\section{計画の必要性}\label{sec:planning}

計画の本質は、意志の力が充実しているときに、将来の意志が減退している自分のために
「行動の線路（レール）」を敷いておくことにあります。

\begin{corollary}[無計画な主体は自由ではない]\label{cor:unfree}
計画を持たない主体は、自由意思を行使しているのではなく、実質的に\textbf{フレーム慣性に支配されている}。
「あと5分だけSNSを見よう」という決断は、自由意思ではなく、慣性が下した決断である。
\end{corollary}

% ─────────────────────────────────────────────────────────────────────────────
\section{まとめ}\label{sec:summary}

\begin{insightbox}
一言で言えば、自由意思を持って慣性そのものと真正面から戦おうとしないでください。
そうではなく、適切なタイミングで自由意思を使って慣性の\textbf{方向}を変え、
あとは慣性の力（習慣の力）を使って目的を達成させるのです。
\end{insightbox}
